\section{Introduction}
\label{sec:intro}

Statutes are not free-form prose. A modern penal code is drafted to a tight
internal grammar: numbered sections; enumerated elements; in-line illustration
blocks; cross-cutting general exceptions; penalty clauses with named
combinators (``or both''). Lawyers learn to read this structure as if it were
a typed program, and computer scientists have, since at least Sergot et al.'s
1986 encoding of the British Nationality Act~\cite{sergot1986british}, observed
that the structure is amenable to direct symbolic representation. The persistent
gap is not the recognition that statute looks like code --- it is the practical
absence of a representation expressive enough to encode an entire production
statute without immediately falling back on free-text strings.

This paper argues that the gap is mostly a grammar gap, and presents
\emph{Yuho}, a domain-specific language designed by mass-encoding the entire
Singapore Penal Code 1871. The headline number is engineering rather than
methodological: \statRawSections{} sections, \statEncoded{} encoded,
\statLOnePass{}/\statRawSections{} passing parse-and-build (the L1+L2 tiers of
our coverage harness), and \statLThreePass{} carrying a structured
human-fidelity stamp (L3) at the time of writing. The interesting numbers
sit underneath: every section that the grammar could not natively express
is documented as a numbered \emph{grammar gap} (\textsf{G1}--\textsf{G14}),
and we report which were genuine grammar limitations, which were diagnostic
issues misread as grammar issues, and which we resolved versus deferred.

\paragraph{Why a fresh DSL?} Existing legal-encoding work clusters into three
camps. The logic-programming tradition treats statute as Horn clauses
\cite{sergot1986british,benchcapon2003legal}; the markup tradition treats
statute as structured XML for archival and exchange (Akoma Ntoso
\cite{akomantoso}, LegalRuleML \cite{legalruleml}); and the modern DSL
tradition picks a sub-domain --- French tax code (Catala~\cite{catala2021}),
contracts and regulations (lam4~\cite{lam4}), generic statute compilation
(LexScript~\cite{lexscript}), Word-driven authoring (LDOC~\cite{ldoc}) ---
and builds a focused syntax for it. None of these target a common-law
\emph{penal} code at full coverage, and none expose primitives that match
how penal sections are actually drafted: an offence is a typed conjunction
of elements; a penalty is a small algebra over imprisonment / fine / caning /
death; an exception is a prioritised override; an illustration is a
verbatim worked example tied to the parent section; an amendment is a
new \texttt{effective <date>} clause that does not erase the prior text.

\paragraph{Contributions.} We make four contributions.

\begin{enumerate}
\item \textbf{The Yuho grammar.} A statute-shaped syntax in which the
top-level construct is a \texttt{statute N \{ \dots \}} block with first-class
fields for definitions, elements, penalty, illustrations, exceptions, case
law, parties, and amendment lineage. Elements carry a deontic type
(\texttt{actus\_reus} / \texttt{mens\_rea} / \texttt{circumstance} /
\texttt{obligation} / \texttt{prohibition} / \texttt{permission}), an optional
burden qualifier (\texttt{burden prosecution} / \texttt{burden defence})
with a named proof standard, and optional causal (\texttt{causedBy}) and
temporal (\texttt{precedes} / \texttt{during} / \texttt{after}) edges to
other elements. Penalty clauses compose three top-level combinators
(\texttt{cumulative} / \texttt{alternative} / \texttt{or\_both}) plus a
\texttt{when <ident>} conditional and a nested-combinator escape hatch.
Exceptions form a Catala-style priority DAG with optional \texttt{defeats}
edges to the elements they override.

\item \textbf{A complete public encoding.} All \statRawSections{} sections
of the Singapore Penal Code 1871, version-controlled under
\texttt{library/penal\_code/}, with a per-section directory layout
(\texttt{statute.yh}, \texttt{metadata.toml}, \texttt{test\_statute.yh}).
The encoding is canonical: the \texttt{.yh} file is the source of truth;
all rendered formats are generated.

\item \textbf{Tooling.} A tree-sitter parser, an immutable Python AST,
six transpilers (JSON, controlled English, \LaTeX{}, Mermaid, Alloy,
DOCX), a Z3-and-Alloy verification hookup, a fidelity-diagnostic pass
that cross-checks encodings against the canonical scrape, a Language Server
exposing hover / inlay-hint / completion / code-lens / diagnostics, a
Model Context Protocol server for AI-driven encoding workflows, and a
VS Code extension wrapping the Language Server.

\item \textbf{Empirical findings.} A catalogue of fourteen distinct grammar
gaps surfaced during mass-encoding, classified into ten genuine grammar
limitations (all resolved by extending the parser) and four diagnostic-layer
issues (all resolved by adding fidelity linters). We argue this distinction
matters: a gap that looks like a grammar problem may, on inspection, be a
linting problem in disguise --- and conflating the two leads to grammar
bloat with no corresponding gain in fidelity.
\end{enumerate}

\paragraph{Why criminal law and why Singapore?} The Singapore Penal Code 1871
was selected for three pragmatic reasons. First, it is public: the canonical
text is hosted on Singapore Statutes Online (SSO)~\cite{singaporepenalcode}
under permissive terms that allow programmatic ingestion. Second, it is
modestly sized: \statRawSections{} sections is small enough to mass-encode
with manageable effort, and large enough to surface structural patterns that
hand-picked toy examples would miss. Third, it inherits the Anglo-Indian
drafting tradition --- short numbered sections, in-line illustrations,
``or both'' penalty clauses, Chapter IV general exceptions --- which is
shared with the Indian, Pakistani, Bruneian, and (with adjustments)
Malaysian penal codes. A grammar that fits the Singapore statute is
plausibly portable across the wider Penal Code family, though we have not
yet empirically verified that.

\paragraph{What this paper is not.} It is not a claim that the encoded
sections constitute legal advice, that the L3 fidelity stamp replaces
external counsel review, or that Yuho is the right tool for every
legal-encoding problem (regulatory corpora, contract drafting, treaty
interpretation, and judgment summarisation each have grammars that Yuho
does not natively model). It is a claim that the gap between a drafted
penal section and an executable artefact is smaller than the existing
literature suggests --- and that closing that gap, for one statute family
at full coverage, is mostly an engineering exercise once the grammar is
right.

\paragraph{Roadmap.} \S\ref{sec:background} surveys the prior art in
statute encoding and positions Yuho against the three traditions.
\S\ref{sec:design} presents the grammar through worked examples.
\S\ref{sec:implementation} describes the toolchain end-to-end.
\S\ref{sec:evaluation} reports empirical results from mass-encoding the
Penal Code, including the gap-trigger frequency analysis and the fidelity
diagnostic hit rate. \S\ref{sec:related} compares Yuho against Catala,
lam4, LexScript, Akoma Ntoso, LegalRuleML, and LDOC. \S\ref{sec:limitations}
discusses what we cannot do, and \S\ref{sec:conclusion} concludes.
